\appendix
\chapter{Complete source-level typing rules}\label{app:source-rules}

This appendix lists the complete set of typing rules for the Veritas source language. Rules marked with $\star$ are also presented in Chapter~\ref{ch:impl}; they are repeated here for completeness.

\section*{Subsumption}

$$\frac{\phi; \Delta; \Gamma \vdash e \Rightarrow (\phi'; \Delta'; \tau_e) \quad \phi' \vdash \tau_e <: \tau}{\phi; \Delta; \Gamma \vdash e \Leftarrow \tau : (\phi'; \Delta')} \quad (\textsc{T-Sub})^\star$$

\section*{Subtyping}

$$\frac{}{\phi \vdash \tau <: \tau} \quad (\textsc{Sub-Refl})$$

$$\frac{}{\phi \vdash \textsf{int}(N) <: \textsf{int}} \quad (\textsc{Sub-Singleton-Int})$$

$$\frac{}{\phi \vdash \{x: \textsf{int} \mid P\} <: \textsf{int}} \quad (\textsc{Sub-Refine-Int})$$

$$\frac{\phi \vdash P[N/x]}{\phi \vdash \textsf{int}(N) <: \{x: \textsf{int} \mid P\}} \quad (\textsc{Sub-Singleton-Refine})^\star$$

$$\frac{\phi \wedge P_1 \vdash P_2}{\phi \vdash \{x: \textsf{int} \mid P_1\} <: \{x: \textsf{int} \mid P_2\}} \quad (\textsc{Sub-Refine-Refine})^\star$$

$$\frac{\phi \vdash \tau_1 <: \tau_2}{\phi \vdash [\tau_1; N] <: [\tau_2; N]} \quad (\textsc{Sub-Array})$$

$$\frac{\phi \vdash \tau <: \sigma}{\phi \vdash \textsf{master}(\tau) <: \sigma} \quad (\textsc{Sub-Master})$$

\section*{Literals and variables}

$$\frac{}{\phi; \Delta; \Gamma \vdash n \Rightarrow (\phi; \Delta; \textsf{int}(n))} \quad (\textsc{T-Int})$$

$$\frac{}{\phi; \Delta; \Gamma \vdash \textsf{true} \Rightarrow (\phi; \Delta; \textsf{bool})} \quad (\textsc{T-True})
  \qquad
  \frac{}{\phi; \Delta; \Gamma \vdash \textsf{false} \Rightarrow (\phi; \Delta; \textsf{bool})} \quad (\textsc{T-False})$$

$$\frac{(x: \tau) \in \Gamma}{\phi; \Delta; \Gamma \vdash x \Rightarrow (\phi; \Delta; \tau)} \quad (\textsc{T-Var-Imm})
  \qquad
  \frac{(x: \tau_c) \in \Delta}{\phi; \Delta; \Gamma \vdash x \Rightarrow (\phi; \Delta; \tau_c)} \quad (\textsc{T-Var-Mut})$$

\section*{Arithmetic}

$$
  \frac{
    \begin{array}{c}
      \phi_0; \Delta_0; \Gamma \vdash e_L \Rightarrow (\phi_1; \Delta_1; \tau_L) \quad \phi_1 \vdash \tau_L <: \textsf{int} \\
      \phi_1; \Delta_1; \Gamma \vdash e_R \Rightarrow (\phi_2; \Delta_2; \tau_R) \quad \phi_2 \vdash \tau_R <: \textsf{int} \\
      \tau_\textsf{res} = \textsf{join-op}(\textsf{aop}, \tau_L, \tau_R)
    \end{array}
  }{
    \phi_0; \Delta_0; \Gamma \vdash e_L \textsf{ aop } e_R \Rightarrow (\phi_2; \Delta_2; \tau_\textsf{res})
  } \quad (\textsc{T-Op})^\star
$$

\section*{Comparison and logical operators}

$$
  \frac{
    \begin{array}{c}
      \phi_0; \Delta_0; \Gamma \vdash e_L \Rightarrow (\phi_1; \Delta_1; \tau_L) \quad \phi_1 \vdash \tau_L <: \textsf{int} \\
      \phi_1; \Delta_1; \Gamma \vdash e_R \Rightarrow (\phi_2; \Delta_2; \tau_R) \quad \phi_2 \vdash \tau_R <: \textsf{int}
    \end{array}
  }{
    \phi_0; \Delta_0; \Gamma \vdash e_L \textsf{ cop } e_R \Rightarrow (\phi_2; \Delta_2; \textsf{bool})
  } \quad (\textsc{T-Cmp})
$$

$$
  \frac{
    \begin{array}{c}
      \phi_0; \Delta_0; \Gamma \vdash e_L \Rightarrow (\phi_1; \Delta_1; \tau_L) \quad \phi_1 \vdash \tau_L <: \textsf{bool} \\
      \phi_1; \Delta_1; \Gamma \vdash e_R \Rightarrow (\phi_2; \Delta_2; \tau_R) \quad \phi_2 \vdash \tau_R <: \textsf{bool}
    \end{array}
  }{
    \phi_0; \Delta_0; \Gamma \vdash e_L \textsf{ lop } e_R \Rightarrow (\phi_2; \Delta_2; \textsf{bool})
  } \quad (\textsc{T-Lop})
$$

\section*{Array creation}

$$
  \frac{
    \begin{array}{c}
      \phi_0; \Delta_0; \Gamma \vdash e_\textsf{val} \Rightarrow (\phi_1; \Delta_1; \tau_\textsf{val})   \\
      \phi_1; \Delta_1; \Gamma \vdash e_\textsf{size} \Rightarrow (\phi_2; \Delta_2; \tau_\textsf{size}) \\
      \phi_2 \vdash \tau_\textsf{size} <: \textsf{int} \quad N = \textsf{val}(\tau_\textsf{size})
    \end{array}
  }{
    \phi_0; \Delta_0; \Gamma \vdash [e_\textsf{val}; e_\textsf{size}] \Rightarrow (\phi_2; \Delta_2; [\tau_\textsf{val}; N])
  } \quad (\textsc{T-Array-Create})
$$

\section*{Array indexing}

$$
  \frac{
    \begin{array}{c}
      (x: [\tau_\textsf{elem}; N]) \in (\Gamma \cup \Delta_0)                                                               \\
      \phi_0; \Delta_0; \Gamma \vdash e_i \Rightarrow (\phi_1; \Delta_1; \tau_i) \quad \phi_1 \vdash \tau_i <: \textsf{int} \\
      \phi_1 \wedge \textsf{prop}(\tau_i) \vdash 0 \le \textsf{val}(\tau_i) < N \quad \text{(Bounds proof)}
    \end{array}
  }{
    \phi_0; \Delta_0; \Gamma \vdash x[e_i] \Rightarrow (\phi_1; \Delta_1; \tau_\textsf{elem})
  } \quad (\textsc{T-Array-Idx})^\star
$$

\section*{Conditional expressions}

$$
  \frac{
    \begin{array}{c}
      \phi_0; \Delta_0; \Gamma \vdash e_\textsf{cond} \Rightarrow (\phi_1; \Delta_1; \tau_c) \quad \phi_1 \vdash \tau_c <: \textsf{bool} \\
      \phi_1 \wedge \textsf{prop}(\tau_c); \Delta_1; \Gamma \vdash B \Rightarrow (\phi_2; \Delta_2; \tau_2)                              \\
      \phi_1 \wedge \neg \textsf{prop}(\tau_c); \Delta_1; \Gamma \vdash B' \Rightarrow (\phi_3; \Delta_3; \tau_3)                        \\
      (\phi_\textsf{join}, \Delta_\textsf{join}, \tau_\textsf{join}) = \textsf{join}(\phi_2, \Delta_2, \tau_2, \phi_3, \Delta_3, \tau_3)
    \end{array}
  }{
    \phi_0; \Delta_0; \Gamma \vdash (\textsf{if } e_\textsf{cond}\; B \textsf{ else } B') \Rightarrow (\phi_\textsf{join}; \Delta_\textsf{join}; \tau_\textsf{join})
  } \quad (\textsc{T-If-Expr})^\star
$$

\section*{Function calls}

$$
  \frac{
    \begin{array}{c}
      \Sigma_F(f) = ( (x_i:\tau_i)_{i=1}^k \rightarrow \tau_\textsf{ret}\; \textsf{requires } P_\textsf{req} ) \\
      \phi_0; \Delta_0; \Gamma \vdash \vec{e} \Rightarrow (\phi_k; \Delta_k; \vec{\sigma})                     \\
      \forall i \in 1..k,\; \phi_k \vdash \sigma_i <: \tau_i                                                   \\
      \phi_k \wedge \textsf{prop}(\vec{\sigma}) \vdash P_\textsf{req}[\textsf{val}(\vec{\sigma}) / \vec{x}] \quad \text{(Precondition proof)}
    \end{array}
  }{
    \phi_0; \Delta_0; \Gamma \vdash f(\vec{e}) \Rightarrow (\phi_k; \Delta_k; \tau_\textsf{ret})
  } \quad (\textsc{T-Call})^\star
$$

\section*{Blocks and bindings}

$$
  \frac{}{\phi; \Delta; \Gamma \vdash \cdot : (\phi; \Delta; \Gamma)} \quad (\textsc{T-Block-Empty})
$$

$$
  \frac{
    \phi_0; \Delta_0; \Gamma_0 \vdash S : (\phi_1; \Delta_1; \Gamma_1) \quad \phi_1; \Delta_1; \Gamma_1 \vdash B : (\phi_2; \Delta_2; \Gamma_2)
  }{
    \phi_0; \Delta_0; \Gamma_0 \vdash (S\; B) : (\phi_2; \Delta_2; \Gamma_2)
  } \quad (\textsc{T-Block-Seq})
$$

$$\frac{\phi_0; \Delta_0; \Gamma \vdash e \Leftarrow \tau : (\phi_1; \Delta_1)}{\phi_0; \Delta_0; \Gamma \vdash (\textsf{let } x:\tau = e;) : (\phi_1; \Delta_1; \Gamma, x:\tau)} \quad (\textsc{T-Let})$$

$$\frac{\phi_0; \Delta_0; \Gamma \vdash e \Leftarrow \tau : (\phi_1; \Delta_1)}{\phi_0; \Delta_0; \Gamma \vdash (\textsf{let mut } x:\tau = e;) : (\phi_1; (\Delta_1, x:\tau); \Gamma)} \quad (\textsc{T-Let-Mut})$$

\section*{Array assignment}

$$
  \frac{
    \begin{array}{c}
      (x: [\tau_\textsf{elem}; N]) \in \Delta_0                                                                                 \\
      \phi_0; \Delta_0; \Gamma \vdash e_i \Rightarrow (\phi_1; \Delta_1; \tau_i) \quad \phi_1 \vdash \tau_i <: \textsf{int}     \\
      \phi_1; \Delta_1; \Gamma \vdash e \Rightarrow (\phi_2; \Delta_2; \tau_v) \quad \phi_2 \vdash \tau_v <: \tau_\textsf{elem} \\
      \phi_2 \wedge \textsf{prop}(\tau_i) \vdash 0 \le \textsf{val}(\tau_i) < N \quad \text{(Bounds proof)}
    \end{array}
  }{
    \phi_0; \Delta_0; \Gamma \vdash (x[e_i] = e;) : (\phi_2; \Delta_2; \Gamma)
  } \quad (\textsc{T-Array-Assign})
$$

\section*{Assignment}

$$
  \frac{
    \begin{array}{c}
      (x: \tau_m) \in \Delta_0                                                 \\
      \phi_0; \Delta_0; \Gamma \vdash e \Rightarrow (\phi_1; \Delta_1; \tau_e) \\
      \phi_1 \vdash \tau_e <: \tau_m
    \end{array}
  }{
    \phi_0; \Delta_0; \Gamma \vdash (x = e;) : (\phi_1; \Delta_1[x \mapsto (\tau_e, \tau_m)]; \Gamma)
  } \quad (\textsc{T-Assign})^\star
$$

\section*{For loop}

$$
  \frac{
    \begin{array}{c}
      \phi_0 \wedge P_I[e_\textsf{start}/i] \quad \text{(Invariant holds on entry)}                                                                   \\
      \phi_0 \wedge P_I \wedge e_\textsf{start} \le i < e_\textsf{end};\; \Delta_0;\; (\Gamma, i: \textsf{int}) \vdash B : (\phi_1; \Delta_1; \Gamma) \\
      \phi_1 \vdash P_I[(i+1)/i] \quad \text{(Invariant preserved by body)}
    \end{array}
  }{
    \phi_0; \Delta_0; \Gamma \vdash (\textsf{for } i \textsf{ in } e_\textsf{start}..e_\textsf{end}\;\{B\}) : (\phi_0 \wedge P_I; \Delta_0; \Gamma)
  } \quad (\textsc{T-For})^\star
$$

\noindent where $P_I$ is the invariant proposition (or $\top$ if no invariant is specified).

\section*{Return}

$$
  \frac{
    \begin{array}{c}
      \phi_0; \Delta_0; \Gamma \vdash e \Rightarrow (\phi_1; \Delta_1; \tau_e) \quad \phi_1 \vdash \tau_e <: \tau_\textsf{ret} \\
      \phi_1 \vdash P_\textsf{ens}[e/\mathtt{result}] \quad \text{(Postcondition proof)}
    \end{array}
  }{
    \phi_0; \Delta_0; \Gamma \vdash (\textsf{return } e;) : (\phi_1; \Delta_1; \Gamma)
  } \quad (\textsc{T-Return})^\star
$$

\section*{If-statement and while loop}

$$
  \frac{
    \begin{array}{c}
      \phi_0; \Delta_0; \Gamma \vdash e_\textsf{cond} \Rightarrow (\phi_1; \Delta_1; \tau_c) \quad \phi_1 \vdash \tau_c <: \textsf{bool} \\
      \phi_1 \wedge \textsf{prop}(\tau_c); \Delta_1; \Gamma \vdash B : (\phi_2; \Delta_2; \Gamma)                                        \\
      \phi_1 \wedge \neg \textsf{prop}(\tau_c); \Delta_1; \Gamma \vdash B' : (\phi_3; \Delta_3; \Gamma)                                  \\
      (\phi_\textsf{join}, \Delta_\textsf{join}) = \textsf{join}(\phi_2, \Delta_2, \phi_3, \Delta_3)
    \end{array}
  }{
    \phi_0; \Delta_0; \Gamma \vdash (\textsf{if } e_\textsf{cond}\; B \textsf{ else } B') : (\phi_\textsf{join}; \Delta_\textsf{join}; \Gamma)
  } \quad (\textsc{T-If-Stmt})
$$

$$
  \frac{
    \begin{array}{c}
      \phi_0; \Delta_0; \Gamma \vdash e_\textsf{cond} \Rightarrow (\phi_1; \Delta_1; \tau_c) \quad \phi_1 \vdash \tau_c <: \textsf{bool} \\
      \phi_0 \vdash \phi_I \quad \text{(Invariant holds on entry)}                                                                       \\
      \phi_I \wedge \textsf{prop}(\tau_c);\; \Delta_0;\; \Gamma \vdash B : (\phi_2; \Delta_2; \Gamma)                                    \\
      \phi_2; \Delta_2 \vdash \phi_I \quad \text{(Invariant preserved by body)}
    \end{array}
  }{
    \phi_0; \Delta_0; \Gamma \vdash (\textsf{while } e_\textsf{cond}\;\{B\}) : (\phi_I \wedge \neg\textsf{prop}(\tau_c); \Delta_0; \Gamma)
  } \quad (\textsc{T-While})
$$

\chapter{Complete DTAL typing rules}\label{app:dtal-rules}

This appendix lists the complete set of typing rules implemented by the DTAL verifier. Each rule has the form $\Phi; \Gamma \vdash \textsf{instr} \Rightarrow \Gamma'$, where $\Phi$ is the constraint context, $\Gamma$ maps registers to types, and $\Gamma'$ is the updated register map. We write $\textsf{idx}(\tau)$ for the index expression extracted from a type (e.g., $\textsf{idx}(\textsf{int}(e)) = e$). Rules marked with $\star$ are also presented in Chapter~\ref{ch:impl}.

\section*{Data movement}

$$
  \frac{}{\Phi; \Gamma \vdash \textsf{mov } r_d, n \Rightarrow \Gamma[r_d \mapsto \textsf{int}(n)]} \quad (\textsc{D-MovImm})^\star
$$

$$
  \frac{\Gamma(r_s) = \tau}{\Phi; \Gamma \vdash \textsf{mov } r_d, r_s \Rightarrow \Gamma[r_d \mapsto \tau]} \quad (\textsc{D-MovReg})
$$

$$
  \frac{
    \begin{array}{c}
      \Gamma(r_b) = [\tau; N] \quad \Phi \vdash 0 \le \textsf{idx}(\Gamma(r_i)) < N
    \end{array}
  }{\Phi; \Gamma \vdash \textsf{load } r_d, [r_b, r_i] \Rightarrow \Gamma[r_d \mapsto \tau]} \quad (\textsc{D-Load})^\star
$$

$$
  \frac{
    \begin{array}{c}
      \Gamma(r_b) = [\tau; N] \quad \Phi \vdash 0 \le \textsf{idx}(\Gamma(r_i)) < N \\
      v' = v + 1 \quad \text{(fresh array version)}
    \end{array}
  }{\Phi; \Gamma \vdash \textsf{store } [r_b, r_i], r_s \Rightarrow \Gamma} \quad (\textsc{D-Store})
$$

\noindent After a store, the verifier emits two axioms into $\Phi$: a \textit{write axiom} ($\textsf{select}(\mathit{arr}_{v'}, e_i) = e_s$) and a \textit{frame axiom} ($\forall k.\; k \ne e_i \Rightarrow \textsf{select}(\mathit{arr}_{v'}, k) = \textsf{select}(\mathit{arr}_v, k)$), where $v$ and $v'$ are the old and new array versions.

\section*{Arithmetic}

$$
  \frac{
    \Gamma(r_1) = \textsf{int}(e_1) \quad \Gamma(r_2) = \textsf{int}(e_2)
  }{\Phi; \Gamma \vdash \textsf{binop}_\oplus\; r_d, r_1, r_2 \Rightarrow \Gamma[r_d \mapsto \textsf{int}(e_1 \oplus e_2)]} \quad (\textsc{D-BinOp-Singleton})^\star
$$

$$
  \frac{
    \Gamma(r_1) = \{a\!:\!\textsf{int} \mid \varphi_1\} \quad \Gamma(r_2) = \{b\!:\!\textsf{int} \mid \varphi_2\}
  }{\Phi; \Gamma \vdash \textsf{binop}_\oplus\; r_d, r_1, r_2 \Rightarrow \Gamma[r_d \mapsto \exists c.\; \textsf{int}(c) \;\mathbf{where}\; \exists a,b.\; \varphi_1 \wedge \varphi_2 \wedge c = a \oplus b]} \quad (\textsc{D-BinOp-Refined})^\star
$$

$$
  \frac{
    \Gamma(r_s) = \textsf{int}(e)
  }{\Phi; \Gamma \vdash \textsf{addimm } r_d, r_s, n \Rightarrow \Gamma[r_d \mapsto \textsf{int}(e + n)]} \quad (\textsc{D-AddImm})
$$

$$
  \frac{
    \Gamma(r_s) = \tau \quad \tau <: \textsf{bool}
  }{\Phi; \Gamma \vdash \textsf{not } r_d, r_s \Rightarrow \Gamma[r_d \mapsto \textsf{bool}]} \quad (\textsc{D-Not})
$$

\section*{Comparison and branching}

$$
  \frac{
    \Gamma(r_1) = \tau_1 \quad \Gamma(r_2) = \tau_2
  }{\Phi; \Gamma \vdash \textsf{cmp } r_1, r_2 \Rightarrow \Gamma \quad \text{(records operands } \textsf{idx}(\tau_1), \textsf{idx}(\tau_2)\text{)}} \quad (\textsc{D-Cmp})
$$

$$
  \frac{
    \Gamma(r_1) = \tau_1
  }{\Phi; \Gamma \vdash \textsf{cmp } r_1, n \Rightarrow \Gamma \quad \text{(records operands } \textsf{idx}(\tau_1), n\text{)}} \quad (\textsc{D-CmpImm})
$$

$$
  \frac{}{\Phi; \Gamma \vdash \textsf{setcc } r_d, \textit{cond} \Rightarrow \Gamma[r_d \mapsto \textsf{bool}]} \quad (\textsc{D-SetCC})
$$

$$
  \frac{
    \text{last cmp operands: } e_1, e_2
  }{\Phi; \Gamma \vdash \textsf{branch } \textit{cond}\; L \Rightarrow \Gamma \quad \text{with } \Phi' = \Phi \wedge \neg(e_1 \textit{ cond } e_2) \text{ (fall-through)}} \quad (\textsc{D-Branch})^\star
$$

\noindent The taken edge to $L$ carries $\Phi \wedge (e_1 \textit{ cond } e_2)$.

\section*{Stack and allocation}

$$
  \frac{\Gamma(r_s) = \tau}{\Phi; \Gamma \vdash \textsf{push } r_s \Rightarrow \Gamma \quad \text{(push } \tau \text{ onto typed stack)}} \quad (\textsc{D-Push})
$$

$$
  \frac{S = \tau :: S'}{\Phi; \Gamma \vdash \textsf{pop } r_d \Rightarrow \Gamma[r_d \mapsto \tau] \quad \text{(pop from typed stack)}} \quad (\textsc{D-Pop})
$$

$$
  \frac{}{\Phi; \Gamma \vdash \textsf{alloca } r_d, n \Rightarrow \Gamma[r_d \mapsto [\textsf{int}; n]]} \quad (\textsc{D-Alloca})
$$

\noindent After allocation, the verifier adds a zero-initialisation axiom: $\forall k \in [0, n).\; \textsf{select}(\mathit{arr}_0, k) = 0$.

\section*{Annotations}

$$
  \frac{
    \Gamma(r) = \tau_\textsf{derived} \quad \Phi \vdash \tau_\textsf{derived} <: \tau_\textsf{ann}
  }{\Phi; \Gamma \vdash \textsf{type } r: \tau_\textsf{ann}} \quad (\textsc{D-TypeAnn})
$$

\noindent The verifier derives the register's type independently and checks it is compatible with the compiler's annotation. If the annotation is more precise (e.g., a singleton vs.\ the derived existential), the verifier checks that the derived type implies the annotation under $\Phi$.

$$
  \frac{
    \Phi \vdash C
  }{\Phi; \Gamma \vdash \textsf{assert } C \Rightarrow \Gamma \quad \text{(add } C \text{ to } \Phi\text{)}} \quad (\textsc{D-Assert})
$$

\noindent The verifier checks that the constraint $C$ is provable from the current $\Phi$ before adding it.

\section*{Function calls}

$$
  \frac{
    \begin{array}{c}
      \Sigma(f) = (\vec{x}:\vec{\tau}) \rightarrow \tau_\textsf{ret} \;\textsf{requires } P_\textsf{req} \;\textsf{ensures } P_\textsf{ens} \\
      \Phi \vdash P_\textsf{req}[\overline{r_i / x_i}]
    \end{array}
  }{\Phi; \Gamma \vdash \textsf{call } f \Rightarrow \Gamma[r_0 \mapsto \tau_\textsf{ret}] \quad \text{(add } P_\textsf{ens}[\overline{r_i / x_i}] \text{ to } \Phi\text{)}} \quad (\textsc{D-Call})
$$

\noindent The precondition is checked against $\Phi$ after substituting physical parameter registers for formal parameters. The postcondition (with \texttt{result} replaced by $r_0$) is assumed in $\Phi$ after the call.

\section*{Physical instructions (post-regalloc)}

The following instructions appear only after register allocation and are verified for basic consistency:

\begin{itemize}
  \item \textbf{D-Cqo}: Sign-extends \texttt{rax} into \texttt{rdx:rax}. The type of \texttt{rdx} is set to the type of \texttt{rax}.
  \item \textbf{D-Idiv}: Divides \texttt{rdx:rax} by the source register. Both \texttt{rax} (quotient) and \texttt{rdx} (remainder) receive type $\textsf{int}$.
  \item \textbf{D-SpillStore}: Records the type stored at a stack offset.
  \item \textbf{D-SpillLoad}: Checks that the loaded type matches what was previously stored at that offset.
  \item \textbf{D-Prologue / D-Epilogue}: Structural markers verified at the function level.
\end{itemize}

\chapter{Evaluation tables}

\begin{table}[H]
  \centering
  \footnotesize
  \setlength{\tabcolsep}{3.6pt}
  \renewcommand{\arraystretch}{0.98}
  \begin{tabular}{p{0.33\textwidth} c c p{0.4\textwidth}}
    \hline
    \textbf{Feature}                           & \textbf{Xi--Harper} & \textbf{Veritas} & \textbf{Notes}                                                                                                               \\
    \hline
    Independent low-level verification         & Yes                 & Yes              & Core architectural idea shared by both systems.                                                                              \\
    Dependent block / label state types        & Yes                 & Yes              & Xi--Harper uses typed labels; Veritas uses \texttt{.entry} plus \texttt{.assume}.                                            \\
    Singleton integer types \texttt{int(x)}    & Yes                 & Yes              & Present in both; Veritas derives symbolic index expressions similarly.                                                       \\
    Size-indexed array types                   & Yes                 & Yes              & Xi--Harper uses \texttt{int array(m)}; Veritas uses array types indexed by \texttt{IndexExpr}.                               \\
    Array-bound-check elimination support      & Yes                 & Yes              & Central purpose in Xi--Harper; preserved in Veritas's array reasoning.                                                       \\
    Typed stack discipline                     & Yes                 & Yes              & Present in both, though exposed differently in the concrete syntax.                                                          \\
    Existential dependent types                & Yes                 & Partial          & Xi--Harper uses them centrally; Veritas currently supports existential integers, but not the broader original style.         \\
    Tuple / product heap objects               & Yes                 & No               & Xi--Harper has \texttt{prod(...)} and \texttt{newtuple}; Veritas does not currently expose general tuple allocation in DTAL. \\
    Sum / tag reasoning                        & Yes                 & No               & Xi--Harper extends DTAL for sums and tag-check elimination; Veritas does not currently have this DTAL feature.               \\
    Explicit array allocation primitives       & Yes                 & No               & Xi--Harper has \texttt{newarray} and \texttt{arraysize}; Veritas uses a different allocation story.                          \\
    Ownership / borrow operations              & No                  & Yes              & Veritas adds \texttt{move\_owned}, \texttt{alias\_borrow}, \texttt{borrow\_mut}, \texttt{borrow\_end}, \texttt{drop\_owned}. \\
    Mutable vs shared alias discipline         & No                  & Yes              & Not present in Xi--Harper; central extension in Veritas.                                                                     \\
    Physical post-regalloc verification target & No                  & Yes              & Veritas verifies physically allocated DTAL after register allocation.                                                        \\
    ABI / calling-convention instructions      & No                  & Yes              & Veritas includes \texttt{spill\_store}, \texttt{spill\_load}, \texttt{.prologue}, \texttt{.epilogue}.                        \\
    Machine-level division expansion           & No                  & Yes              & Veritas includes physical instructions such as \texttt{cqo} and \texttt{idiv}.                                               \\
    \hline
  \end{tabular}
  \caption{Feature comparison between Xi and Harper's DTAL and the DTAL implemented in Veritas.}
  \label{tab:dtal-xi-harper-comparison}
\end{table}
\begin{table}[h]
  \centering
  \small
  \begin{tabular}{l r}
    \hline
    \textbf{Error class}           & \textbf{Total} \\
    \hline
    Type error                     & 7              \\
    Name resolution                & 2              \\
    Control flow                   & 3              \\
    Array error                    & 2              \\
    Call error                     & 2              \\
    Contract violation             & 2              \\
    Quantifier error               & 3              \\
    Invariant failure              & 2              \\
    Safety proof failure           & 2              \\
    Arithmetic safety              & 2              \\
    Overflow safety                & 5              \\
    Other grouped frontend classes & 3              \\
    \hline
  \end{tabular}
  \caption{Rejection of invalid programs in the curated negative suite.}
  \label{tab:negative-suite-summary}
\end{table}

\begin{table}[h]
  \centering
  \small
  \begin{tabular}{l r r r}
    \hline
    \textbf{Kernel}              & \textbf{Veritas (s)} & \textbf{GCC \texttt{-O2} (s)} & \textbf{Ratio} \\
    \hline
    \texttt{floyd\_warshall\_500} & 0.601975           & 0.091029                      & 6.613          \\
    \texttt{jacobi\_1d\_400\_100} & 0.000796          & 0.001396                      & 0.570          \\
    \texttt{jacobi\_2d\_250\_100} & 0.098254          & 0.019011                      & 5.168          \\
    \texttt{seidel\_2d\_400\_100} & 0.316906          & 0.090827                      & 3.489          \\
    \texttt{mvt\_400}            & 0.004776            & 0.003371                      & 1.417          \\
    \texttt{atax\_390\_410}      & 0.005116            & 0.003447                      & 1.484          \\
    \texttt{gesummv\_250}        & 0.002868            & 0.004430                      & 0.647          \\
    \hline
  \end{tabular}
  \caption{Selected PolyBench-style runtime results against GCC
    \texttt{-O2}. All kernels in the evaluated set passed checksum-based
    semantic validation before timing results were accepted.}
  \label{tab:polybench-practicality}
\end{table}

\begin{table}[h]
  \centering
  \small
  \begin{tabular}{l r r r r}
    \hline
    \textbf{Task}                  & \textbf{Veritas} & \textbf{Dafny} & \textbf{Liquid Haskell} & \textbf{Verus} \\
    \hline
    \texttt{safe\_division}        & 0.167            & 0.125          & 0.400                   & 0.100          \\
    \texttt{bounded\_read\_offset} & 0.111            & 0.267          & 0.286                   & 0.182          \\
    \texttt{fill\_with\_ones}      & 0.154            & 0.267          & 0.364                   & 0.250          \\
    \texttt{binary\_search}        & 0.077            & 0.229          & 0.200                   & 0.222          \\
    \texttt{sorted\_head}          & 0.125            & 0.263          & 0.417                   & 0.235          \\
    \texttt{safe\_midpoint}        & 0.143            & 0.250          & 0.400                   & 0.200          \\
    \hline
  \end{tabular}
  \caption{Annotation-to-code ratios for the normalised cross-system comparison suite.}
  \label{tab:cross-system-annotation-burden}
\end{table}

\begin{table}[h]
  \centering
  \small
  \begin{tabularx}{\linewidth}{
      >{\ttfamily}l
      >{\centering\arraybackslash}X
      >{\centering\arraybackslash}X
      >{\centering\arraybackslash}X
      >{\centering\arraybackslash}X
      >{\centering\arraybackslash}X
    }
    \hline
    \textbf{Task}                       &
    \textbf{Veritas compile-only (s)}   &
    \textbf{Veritas compile+verify (s)} &
    \textbf{Dafny (s)}                  &
    \textbf{Liquid Haskell (s)}         &
    \textbf{Verus (s)}                                                      \\
    \hline

    safe\_division                      &
    0.0100                              & 0.0138 & 0.9259 & 1.1794 & 0.4484 \\

    bounded\_read\_offset               &
    0.0171                              & 0.0204 & 1.0118 & 1.2807 & 0.4604 \\

    fill\_with\_ones                    &
    0.0336                              & 0.0398 & 1.0015 & 1.2213 & 0.4924 \\

    binary\_search                      &
    0.1416                              & 0.1780 & 1.2629 & 1.3219 & 0.5285 \\

    sorted\_head                        &
    0.0147                              & 0.0174 & 1.0054 & 1.1942 & 0.5068 \\

    safe\_midpoint                      &
    0.0116                              & 0.0149 & 0.8856 & 1.1624 & 0.4939 \\
    \hline
  \end{tabularx}

  \caption{Exact mean wall-clock timings for the normalised six-task
    cross-system comparison suite.}
  \label{tab:comparison-suite-wallclock}
\end{table}
